%! AP Statistics — Unit 3, Lesson 3.4 — Activity KEY
\documentclass[10pt]{article}
\usepackage{apstats-article}
\usepackage{apstats-key}

\begin{document}

\pageheader{Unit 3, Lesson 3.4}{Group Activity KEY: Bias Investigators}
\namepartnerperiod

\vspace{0.1in}

\textbf{Survey 1:} Online sports poll.
\begin{practicebox}
\small
\begin{enumerate}[leftmargin=1.5em, itemsep=5pt]
  \item \ans{Voluntary response bias}
  \item \ans{Highly opinionated sports fans who visit the website and have strong preferences.}
  \item \ans{The results over-represent the opinions of engaged fans; casual or disinterested people are excluded, so the "winner" may not reflect broader public opinion.}
  \item \ans{Randomly sample basketball fans or the general public using a list-based random sample.}
\end{enumerate}
\end{practicebox}

\vspace{0.08in}

\textbf{Survey 2:} Mail survey with low return rate.
\begin{practicebox}
\small
\begin{enumerate}[leftmargin=1.5em, itemsep=5pt]
  \item \ans{No — the original 500 were randomly selected.}
  \item \ans{Nonresponse bias}
  \item \ans{The 180 respondents are disproportionately elderly; elderly households tend to have lower internet usage, so the estimate of internet usage will be too low.}
  \item \ans{Phone or email follow-up with non-respondents; offer an incentive; use a shorter survey.}
\end{enumerate}
\end{practicebox}

\vspace{0.08in}

\textbf{Survey 3:} Landline-only immigration poll.
\begin{practicebox}
\small
\begin{enumerate}[leftmargin=1.5em, itemsep=5pt]
  \item \ans{Undercoverage bias}
  \item \ans{Young adults and lower-income households, who are more likely to rely solely on cell phones.}
  \item \ans{Young adults and immigrants may hold different views on immigration; excluding them distorts the estimate.}
  \item \ans{The results will over-represent the views of older, higher-income adults who have landlines.}
\end{enumerate}
\end{practicebox}

\vspace{0.08in}

\textbf{Survey 4:} Leading question.
\begin{practicebox}
\small
\begin{enumerate}[leftmargin=1.5em, itemsep=5pt]
  \item \ans{Response bias (question-wording bias)}
  \item \ans{"In order to improve student safety" primes students to agree; "ID-badge policy" is framed as a safety measure.}
  \item \ans{Example: "Do you support the new ID-badge policy?" or "How do you feel about the new ID-badge policy?"}
  \item \ans{Too high — the leading framing pressures students to say yes.}
\end{enumerate}
\end{practicebox}

\vspace{0.08in}

\begin{practicebox}
\small
\textbf{Group Synthesis:}

\renewcommand{\arraystretch}{1.6}
\begin{tabularx}{\linewidth}{@{} >{\bfseries}p{0.25\linewidth} X X @{}}
\textbf{Bias} & \textbf{Cause} & \textbf{Fix} \\
\hline
Voluntary response & \ans{Only motivated volunteers respond.} & \ans{Use a random sample.} \\
Undercoverage & \ans{Part of population excluded from sampling frame.} & \ans{Use a complete, up-to-date sampling frame.} \\
Nonresponse & \ans{Selected individuals refuse or fail to respond.} & \ans{Follow up; use incentives; easier survey.} \\
Response (wording) & \ans{Question or context distorts answers.} & \ans{Use neutral, balanced question wording.} \\
\end{tabularx}
\end{practicebox}

\end{document}
